\section{Introduction}

Multi-agent AI systems that recursively delegate tasks face a fundamental structural vulnerability: the \textbf{autonomous drift problem}. As agentic systems spawn sub-agents that spawn further sub-agents, accountability chains become opaque, execution provenance degrades, and systems lose fidelity to their original intent. Without a formal mechanism to preserve the lineage of decisions, every delegation step introduces compounding uncertainty about which agent performed what action, under whose authority, and toward which objective. This erosion of accountability is not merely a governance concern; it is a correctness concern. When drift accumulates beyond a critical threshold, the system may continue executing in a state that no principal would endorse, with no agent capable of reconstructing the causal chain that produced it.

The field has attempted to address drift through two primary mechanisms. \textbf{Time-to-live} (TTL) constraints limit the operational lifetime of an agent or the number of delegation hops permitted before a workflow terminates. \textbf{Depth limits} cap the maximum nesting depth of recursive spawning. Both approaches impose an external, temporal bound on system behavior. Neither approach, however, addresses the fundamental information-theoretic problem: when a sub-agent spawns a further sub-agent, the system retains no structural record of the parent-child relationship that is both immutable and queryable. TTL expiration does not preserve lineage; it simply terminates. Depth limits do not encode intent propagation; they merely truncate it. As a result, existing systems exhibit a structural blind spot: they can enforce when a system stops, but not what the system did or why.

We argue that the insufficiency of TTL and depth limits stems from a missing primitive: an \textbf{immutable parent pointer chain}. In the BX3 architecture, every agent carries a cryptographically signed, append-only chain of references to its ancestors. A spawned agent does not merely receive instructions; it receives an immutable record of its provenance that it cannot alter, delete, or dissociate from. This chain is structurally impossible to discard because it is embedded in the agent's foundational identity, not stored in a mutable attribute. The agent cannot ``forget'' its parent because its identity is defined, in part, by the chain that produced it. This turns drift from a probabilistic risk into a structural impossibility: an agent cannot drift from intent that is cryptographically embedded in its identity.

We formalize drift as follows. Let an agent $a_i$ be defined by its capability set $C_i$, its delegated authority $A_i$, and its current objective $O_i$. When $a_i$ spawns a sub-agent $a_{i+1}$, drift occurs if and only if the objective propagated to $a_{i+1}$ is not a faithful translation of $O_i$, and no verifiable chain exists to recover $O_i$ from the state of $a_{i+1}$. This definition has two components: \textbf{intent violation} (the objective was modified in propagation) and \textbf{lineage opacity} (no mechanism exists to reconstruct the original intent). Drift is thus not a matter of degree --- a system either has sufficient lineage to reconstruct intent, or it does not. TTL and depth limits address the latter condition (lineage opacity) without addressing the former (intent violation). The immutable parent pointer chain addresses both simultaneously: even if intent were violated in propagation, the chain provides a verifiable reconstruction path back to the delegating principal.

BX3 operationalizes this through the \textbf{Behavior, Execution, and eXchange} architecture. Each agent maintains a provably complete ancestry record called the Parent Pointer Chain (PPC). The PPC is immutable because it is sealed at agent creation via SHA-256 forensic sealing; modification of any ancestor record invalidates the chain's integrity. The chain is queryable at any depth via the Chain Query Interface (CQI), enabling any authorized agent or human overseer to reconstruct the full delegation path. The Training Wheels Protocol further constrains drift by enforcing that all external actions execute through a HITL (human-in-the-loop) queue with full provenance logging, ensuring that even if the PPC provides structural immutability, execution fidelity is enforced at runtime.

The remainder of this paper is organized as follows. Section 2 surveys related work in multi-agent governance, delegation protocols, and provenance tracking. Section 3 formalizes the BX3 architecture, including the PPC structure, the CQI interface, and the Truth Gate mechanism that enforces chain integrity. Section 4 presents the Experimental Evaluation, measuring drift resistance under adversarial propagation scenarios and comparing TTL and depth-limited baselines against BX3. Section 5 discusses limitations and edge cases, including chain traversal performance at extreme depths and the operational overhead of immutable sealing. Section 6 concludes with implications for multi-agent system design and open problems in agentic accountability.